\begin{figure}[H]
\centering
\resizebox{0.96\linewidth}{!}{%
\begin{tikzpicture}[>=Stealth,
  card/.style={draw=black!60, fill=black!3, rounded corners=3pt, minimum width=38mm, minimum height=26mm, text width=34mm, align=center, font=\footnotesize\sffamily, inner sep=3pt},
  result/.style={draw=purple!65!black, fill=purple!5, rounded corners=3pt, minimum width=92mm, minimum height=23mm, text width=86mm, align=left, font=\footnotesize\sffamily, inner sep=4pt},
  mechanism/.style={draw=black!45, fill=white, rounded corners=2pt, minimum width=28mm, minimum height=7mm, align=center, font=\scriptsize\sffamily, inner sep=2pt},
  arrow/.style={->, draw=black!65, line width=.9pt, >={Stealth[length=2.5mm]}},
  source/.style={font=\scriptsize\sffamily, fill=white, inner sep=1.5pt, align=center},
]
  \node[card, draw=blue!60!black, fill=blue!4] (fixed) at (0,0) {\textbf{常规固定翼}\\[2pt]
    \textbf{核心机制：}气动升力面与舵面\\
    \textbf{优势：}巡航效率高\\
    \textbf{局限：}低速控制弱、依赖跑道};
  \node[card, draw=green!55!black, fill=green!4] (quad) at (4.8,0) {\textbf{四旋翼}\\[2pt]
    \textbf{核心机制：}多电机差速控制\\
    \textbf{优势：}悬停与低速操纵强\\
    \textbf{局限：}高速巡航效率低};
  \node[card, draw=orange!70!black, fill=orange!5] (tilt) at (9.6,0) {\textbf{倾转旋翼}\\[2pt]
    \textbf{核心机制：}推力方向连续倾转\\
    \textbf{优势：}兼顾VTOL与巡航\\
    \textbf{局限：}机构重、过渡控制复杂};

  \node[mechanism, draw=blue!60!black] (lift) at (0,-2.55) {固定翼升力面};
  \node[mechanism, draw=green!55!black] (reaction) at (4.8,-2.55) {反扭矩差动};
  \node[mechanism, draw=orange!70!black] (vector) at (9.6,-2.55) {低速矢推控制};
  \draw[arrow, draw=blue!60!black] (fixed.south) -- (lift.north);
  \draw[arrow, draw=green!55!black] (quad.south) -- (reaction.north);
  \draw[arrow, draw=orange!75!black] (tilt.south) -- (vector.north);

  \coordinate (merge) at (4.8,-3.7);
  \draw[arrow, draw=blue!60!black] (lift.south) |- (merge);
  \draw[arrow, draw=green!55!black] (reaction.south) -- (merge);
  \draw[arrow, draw=orange!75!black] (vector.south) |- (merge);
  \fill[black!65] (merge) circle (2pt);
  \node[source, above=3pt] at (merge) {几何融合与四输入主通道};

  \node[result] (ours) at (4.8,-6.45) {\textbf{本文构型：纵列双发正交单轴矢量推力}\\[3pt]
    \textbf{机制组合：}固定翼巡航升力 + 正交单轴矢推 + 反扭矩差动\\
    \textbf{概念潜力：}总推力、两摆角与差速形成结构清晰的四输入主通道\\
    \textbf{方向约定：}前 CW，$\bm h_f\parallel +x_b$；尾 CCW，$\bm h_t\parallel -x_b$；$\delta_f$ 绕 $z_b$，$\delta_t$ 绕 $y_b$（见图~\ref{fig:thrust_mapping}）\\
    \textbf{验证边界：}交叉项、低油门退化、失效状态及巡航效率尚待标定与试验。};
  \draw[arrow, draw=purple!65!black, line width=1.15pt] (merge) -- (ours.north);
\end{tikzpicture}%
}
\caption{构型的机制来源与概念定位：固定翼升力面、反扭矩差动，以及由两个正交单轴摆座产生的矢量控制。图示表达设计关系，不代表已验证的性能排序或VTOL能力。}
\label{fig:config_comparison}
\end{figure}
